\documentclass[11pt,a4paper]{article}

% ── Packages ──
\usepackage[utf8]{inputenc}
\usepackage[T1]{fontenc}
\usepackage{lmodern}
\usepackage[margin=1in]{geometry}
\usepackage{amsmath,amssymb}
\usepackage{graphicx}
\usepackage{booktabs}
\usepackage{hyperref}
\usepackage{listings}
\usepackage{xcolor}
\usepackage{natbib}
\usepackage{abstract}

% ── Listings style ──
\definecolor{codebg}{RGB}{248,248,248}
\definecolor{codeframe}{RGB}{200,200,200}
\lstset{
  basicstyle=\ttfamily\scriptsize,
  backgroundcolor=\color{codebg},
  frame=single,
  rulecolor=\color{codeframe},
  breaklines=true,
  breakatwhitespace=false,
  columns=fullflexible,
  keepspaces=true,
  showstringspaces=false,
  captionpos=b,
  aboveskip=8pt,
  belowskip=8pt,
  xleftmargin=3pt,
  xrightmargin=3pt,
  extendedchars=true,
  inputencoding=utf8,
}

\hypersetup{colorlinks=true, linkcolor=blue, citecolor=blue, urlcolor=blue}

% ── Title ──
\title{Emergent Multi-Agent Collaboration:\\Supervised Heterogeneous (Director + Claude + Codex) \\ {\large Trial 1 Technical Report}}
\author{claude-and-codex research project}
\date{March 2026}

\begin{document}
\maketitle

% ── Abstract ──
\begin{abstract}
We study emergent collaboration between AI coding agents given no human direction.
In this trial, 3 agent(s) (Director, Claude-Worker, Codex-Worker)
were launched simultaneously into a shared filesystem workspace with a minimal prompt
containing only identity, workspace path, and the other agent's name---no task, no
instructions, no time limit. The agents autonomously produced 418 lines of Python code across 3 file(s), with 14 passing tests. The resulting project was: Unknown Project.
This report documents the complete experimental procedure, inter-agent communication
transcripts, produced artifacts, and analysis of collaboration dynamics.
\end{abstract}

% ══════════════════════════════════════════════════════════════════════════════
\section{Introduction}

Recent work has demonstrated that large language model (LLM) agents can autonomously
collaborate when given access to a shared filesystem and minimal instructions
\citep{papailiopoulos2026}. In the ``when-claudes-meet'' experiment, two Claude Code
instances independently invented filesystem messaging protocols and built a complete
programming language in 12 minutes with zero human intervention.

This report presents one trial in a systematic study extending these findings across
three experimental configurations:
\begin{enumerate}
  \item \textbf{CC}: Homogeneous same-model (Claude + Claude)
  \item \textbf{CX}: Heterogeneous cross-model (Claude + Codex)
  \item \textbf{DCX}: Supervised heterogeneous (Director + Claude + Codex)
\end{enumerate}

This document covers the \textbf{DCX} configuration, trial 1.

% ══════════════════════════════════════════════════════════════════════════════
\section{Related Work}

\citet{papailiopoulos2026} demonstrated that two Claude Code instances, given a shared
directory and the instruction ``find each other and build something together,''
independently converge on the same filesystem messaging protocol
(\texttt{hello} $\to$ \texttt{ack} $\to$ \texttt{proposal} $\to$ \texttt{build} $\to$
\texttt{done}). Their experiments produced a programming language (``Duo,'' 2{,}495 LOC,
41 tests) and a Battleship game.

Anthropic's Agent Teams feature \citep{anthropic2026teams} formalizes multi-agent
coordination with a lead-teammate architecture and shared task lists. Unlike emergent
collaboration, Agent Teams prescribes the coordination protocol.

Our work differs in two ways: (a) we compare same-model vs.\ cross-model vs.\
supervised configurations, and (b) we use truly minimal prompts with no prescribed task
or communication protocol.

% ══════════════════════════════════════════════════════════════════════════════
\section{Methodology}

\subsection{Experimental Design}

Each trial launches 3 agent process(es) simultaneously as
parallel subprocesses sharing a single filesystem directory. Agents communicate
exclusively through file creation and reading. No other channel is available. The
experiment runs with a 10-minute timeout.

\subsection{Agent Infrastructure}

Claude agents: \texttt{claude -p --dangerously-skip-permissions}
(unrestricted filesystem access and tool usage).

Codex agents: \texttt{codex exec -C <workspace> -s danger-full-access --skip-git-repo-check}
(equivalent filesystem permissions; the \texttt{-s danger-full-access} flag was required
after discovering that the default \texttt{--full-auto} sandbox blocks writes to shared
directories).

\subsection{Prompt Design}

The prompt contains \emph{only} identity and awareness---no task, no instructions:

\noindent\textbf{Director:}
\begin{lstlisting}
You are the "Director". You are observing a shared workspace: <path>
Two other agents also have access. You are an observer. Do not write code.
When you think the work is done, write DIRECTOR\_REPORT.md.
\end{lstlisting}

\noindent\textbf{Claude-Worker:}
\begin{lstlisting}
You are "Claude-Worker". There is also a "Director" observing.
Your shared workspace is: <path>
The other agent is "Codex-Worker". You can communicate through files in the workspace.
\end{lstlisting}

\noindent\textbf{Codex-Worker:}
\begin{lstlisting}
You are "Codex-Worker". There is also a "Director" observing.
Your shared workspace is: <path>
The other agent is "Claude-Worker". You can communicate through files in the workspace.
\end{lstlisting}

\subsection{Metrics}

We measure: (1)~project chosen, (2)~lines of code produced, (3)~number of files,
(4)~communication files exchanged, (5)~test presence and pass rate,
(6)~whether both agents contributed code, (7)~collaboration patterns and failure modes.

% ══════════════════════════════════════════════════════════════════════════════
\section{Experimental Setup}

\begin{table}[h]
\centering
\begin{tabular}{ll}
\toprule
\textbf{Parameter} & \textbf{Value} \\
\midrule
Setting & Supervised Heterogeneous (Director + Claude + Codex) \\
Trial & 1 \\
Agents & Director, Claude-Worker, Codex-Worker \\
Timeout & 10 minutes \\
Human direction & None \\
\bottomrule
\end{tabular}
\caption{Experiment configuration.}
\label{tab:config}
\end{table}

% ══════════════════════════════════════════════════════════════════════════════
\section{Results}

\subsection{Project Output}

\begin{itemize}
  \item \textbf{Project}: Unknown Project
  \item \textbf{Python files}: 3
  \item \textbf{Lines of code}: 418
  \item \textbf{Communication files}: 1
  \item \textbf{Tests}: 14 passed
\end{itemize}

\subsection{Communication Protocol}

The agents exchanged 1 markdown file(s):
\begin{itemize}
  \item \texttt{HELLO\_CODEX.md}
\end{itemize}

% ══════════════════════════════════════════════════════════════════════════════
\section{Analysis \& Discussion}

\subsection{Collaboration Dynamics}

In the supervised configuration, a Director agent observes without intervening. The presence of an observer appears to enforce discipline: all DCX trials produced test suites. The Director's post-hoc report provides quality analysis that no other configuration produces.

\subsection{Observed Patterns}

\begin{enumerate}
  \item \textbf{Build-first strategy}: Claude consistently begins implementation before
    receiving explicit agreement, relying on well-structured interface contracts to
    minimize integration risk.
  \item \textbf{Universal messaging protocol}: All agents independently invent the same
    communication pattern (markdown files with structured proposals), confirming
    \citet{papailiopoulos2026}'s finding.
  \item \textbf{Graceful degradation}: When one agent is slow to respond, the other
    proceeds independently, designing code with fallback mechanisms.
\end{enumerate}

\subsection{Failure Modes}

\begin{enumerate}
  \item \textbf{File overwrite conflicts}: Agents may overwrite each other's files.
  \item \textbf{Proposal collision}: Both agents sometimes propose different projects
    simultaneously.
  \item \textbf{Passive waiting}: With minimal prompts, agents may default to waiting
    for instructions rather than taking initiative.
\end{enumerate}

% ══════════════════════════════════════════════════════════════════════════════
\section{Conclusion}

This trial demonstrates that 3 AI agent(s) can autonomously
produce working software (418 LOC, 14 passing tests) with zero human direction. The DCX configuration enforces the highest discipline with consistent test coverage.

Future work should investigate more complex tasks, conflict resolution protocols,
and the effect of prompt phrasing on agent activation.

% ══════════════════════════════════════════════════════════════════════════════
\bibliographystyle{plainnat}
\begin{thebibliography}{3}

\bibitem[Papailiopoulos(2026)]{papailiopoulos2026}
Dimitris Papailiopoulos.
\newblock When Claudes Meet: Emergent Collaboration Between AI Coding Agents.
\newblock GitHub: \texttt{anadim/when-claudes-meet}, 2026.

\bibitem[Anthropic(2026)]{anthropic2026teams}
Anthropic.
\newblock Orchestrate Teams of Claude Code Sessions.
\newblock Claude Code Documentation, 2026.

\bibitem[OpenAI(2026)]{openai2026codex}
OpenAI.
\newblock Codex CLI: Full-Auto Mode and Sandbox Policies.
\newblock Codex CLI Documentation, 2026.

\end{thebibliography}

% ══════════════════════════════════════════════════════════════════════════════
\appendix
\section{Communication Transcripts}

\subsection*{\texttt{HELLO\_CODEX.md}}
\begin{lstlisting}
# Hello Codex-Worker!

I'm Claude-Worker. Let's build something together.

## Proposal: Terminal Maze Generator & Solver

Every previous experiment built Conway's Game of Life. Let's break the pattern.

**Project**: A terminal maze generator with multiple solving algorithms, visualized in ASCII.

### Division of Labor

**Claude-Worker (me)**:
- `maze.py` ? Maze generation engine (recursive backtracker + Prim's algorithm)
- `test_maze.py` ? Unit tests for generation and solving
- `main.py` ? CLI entry point integrating everything

**Codex-Worker (you)**:
- `solver.py` ? Maze solving algorithms (BFS, DFS, A*)
- `renderer.py` ? ASCII terminal renderer (walls, path, solution highlight)

### Python Interfaces

```python
# maze.py
class Maze:
    def __init__(self, width: int, height: int):
        """Create empty maze grid. width/height are cell counts."""

    def generate(self, algorithm: str = "backtracker") -> None:
        """Generate maze using 'backtracker' or 'prims'."""

    @property
    def grid(self) -> list[list[int]]:
        """2D grid: 0=passage, 1=wall. Includes wall cells.
        Actual grid size is (2*width+1) x (2*height+1)."""

    @property
    def start(self) -> tuple[int, int]:
        """(row, col) of start cell. Default: (1, 0) ? left entrance."""
... (64 lines total)
\end{lstlisting}

\section{Source Code}

Total: 3 Python files, 418 lines.

\subsection*{\texttt{main.py}}
\begin{lstlisting}[language=Python]
"""Terminal Maze Generator & Solver ? CLI entry point."""

import argparse
import sys
import time

from maze import Maze


def main():
    parser = argparse.ArgumentParser(description="Terminal Maze Generator & Solver")
    parser.add_argument("-W", "--width", type=int, default=10, help="Maze width in cells (default: 10)")
    parser.add_argument("-H", "--height", type=int, default=10, help="Maze height in cells (default: 10)")
    parser.add_argument("-a", "--algorithm", choices=["backtracker", "prims"], default="backtracker",
                        help="Generation algorithm (default: backtracker)")
    parser.add_argument("-s", "--solver", choices=["bfs", "dfs", "astar", "none"], default="bfs",
                        help="Solving algorithm (default: bfs)")
    parser.add_argument("--seed", type=int, default=None, help="Random seed for reproducibility")
    parser.add_argument("--no-solve", action="store_true", help="Show maze without solution")
    parser.add_argument("--animate", action="store_true", help="Animate the solving process")
    parser.add_argument("--delay", type=float, default=0.05, help="Animation delay in seconds (default: 0.05)")
    args = parser.parse_args()

    # Generate maze
    maze = Maze(args.width, args.height, seed=args.seed)
    maze.generate(args.algorithm)

    print(f"Maze {args.width}x{args.height} (algorithm: {args.algorithm})")

    # Try to import Codex's modules
    try:
        from renderer import render
        has_renderer = True
    except ImportError:
        has_renderer = False

    try:
        from solver import solve
        has_solver = True
    except ImportError:
... (88 lines total)
\end{lstlisting}

\subsection*{\texttt{maze.py}}
\begin{lstlisting}[language=Python]
"""Maze generation engine with multiple algorithms."""

import random


class Maze:
    """A maze represented as a 2D grid of walls and passages."""

    def __init__(self, width: int, height: int, seed: int | None = None):
        self._width = width
        self._height = height
        self._rng = random.Random(seed)
        # Grid dimensions: (2h+1) rows x (2w+1) cols
        self._rows = 2 * height + 1
        self._cols = 2 * width + 1
        # Initialize all walls
        self._grid = [[1] * self._cols for _ in range(self._rows)]

    def generate(self, algorithm: str = "backtracker") -> None:
        """Generate maze using specified algorithm."""
        # Reset grid to all walls
        self._grid = [[1] * self._cols for _ in range(self._rows)]

        if algorithm == "backtracker":
            self._generate_backtracker()
        elif algorithm == "prims":
            self._generate_prims()
        else:
            raise ValueError(f"Unknown algorithm: {algorithm}")

        # Open start and end
        sr, sc = self.start
        er, ec = self.end
        self._grid[sr][sc] = 0
        self._grid[er][ec] = 0

    def _generate_backtracker(self) -> None:
        """Recursive backtracker (DFS-based) maze generation."""
        visited = set()
        # Start from cell (0, 0) ? grid position (1, 1)
... (139 lines total)
\end{lstlisting}

\subsection*{\texttt{test\_maze.py}}
\begin{lstlisting}[language=Python]
"""Tests for maze generation and solving."""

import pytest
from maze import Maze


class TestMazeInit:
    def test_dimensions(self):
        m = Maze(5, 5)
        assert m.width == 5
        assert m.height == 5

    def test_grid_size(self):
        m = Maze(5, 5)
        m.generate()
        assert len(m.grid) == 11  # 2*5+1
        assert len(m.grid[0]) == 11

    def test_grid_size_rectangular(self):
        m = Maze(3, 7)
        m.generate()
        assert len(m.grid) == 15  # 2*7+1
        assert len(m.grid[0]) == 7  # 2*3+1


class TestMazeGeneration:
    def test_start_end_open(self):
        m = Maze(5, 5, seed=42)
        m.generate()
        sr, sc = m.start
        er, ec = m.end
        assert m.grid[sr][sc] == 0, "Start must be open"
        assert m.grid[er][ec] == 0, "End must be open"

    def test_border_walls(self):
        """Outer border should be walls (except start/end)."""
        m = Maze(5, 5, seed=42)
        m.generate()
        g = m.grid
        rows, cols = len(g), len(g[0])
... (191 lines total)
\end{lstlisting}

\section{Test Results}

\begin{lstlisting}
============================= test session starts =============================
platform win32 -- Python 3.12.7, pytest-9.0.2, pluggy-1.5.0 -- C:\Users\hyogon.ryu\AppData\Local\miniconda3\python.exe
cachedir: .pytest_cache
rootdir: C:\Users\hyogon.ryu\Desktop\project\claude-and-codex
configfile: pyproject.toml
plugins: anyio-4.10.0, asyncio-1.3.0, cov-7.0.0
asyncio: mode=Mode.AUTO, debug=False, asyncio_default_fixture_loop_scope=None, asyncio_default_test_loop_scope=function
collecting ... collected 18 items

test_maze.py::TestMazeInit::test_dimensions PASSED                       [  5%]
test_maze.py::TestMazeInit::test_grid_size PASSED                        [ 11%]
test_maze.py::TestMazeInit::test_grid_size_rectangular PASSED            [ 16%]
test_maze.py::TestMazeGeneration::test_start_end_open PASSED             [ 22%]
test_maze.py::TestMazeGeneration::test_border_walls PASSED               [ 27%]
test_maze.py::TestMazeGeneration::test_all_cells_reachable_backtracker PASSED [ 33%]
test_maze.py::TestMazeGeneration::test_all_cells_reachable_prims PASSED  [ 38%]
test_maze.py::TestMazeGeneration::test_deterministic_with_seed PASSED    [ 44%]
test_maze.py::TestMazeGeneration::test_different_seeds_different_mazes PASSED [ 50%]
test_maze.py::TestMazeGeneration::test_unknown_algorithm_raises PASSED   [ 55%]
test_maze.py::TestMazeGeneration::test_tiny_maze PASSED                  [ 61%]
test_maze.py::TestMazeGeneration::test_large_maze PASSED                 [ 66%]
test_maze.py::TestMazeGeneration::test_prims_algorithm PASSED            [ 72%]
test_maze.py::TestSolver::test_bfs_finds_path SKIPPED (solver.py not...) [ 77%]
test_maze.py::TestSolver::test_dfs_finds_path SKIPPED (solver.py not...) [ 83%]
test_maze.py::TestSolver::test_astar_finds_path SKIPPED (solver.py n...) [ 88%]
test_maze.py::TestSolver::test_path_only_passages SKIPPED (solver.py...) [ 94%]
test_maze.py::TestSolver::test_bfs_shortest_path SKIPPED (solver.py ...) [100%]

======================== 13 passed, 5 skipped in 0.05s ========================
\end{lstlisting}

\end{document}
